\documentclass{beamer}

% --- THEME + COLORS ---
\usetheme{Madrid} % clean and minimal

% Custom soft orange / cream palette
\definecolor{softorange}{RGB}{240,173,78}
\definecolor{cream}{RGB}{252,247,239}
\definecolor{darkgrey}{RGB}{45,45,45}

\setbeamercolor{structure}{fg=softorange}
\setbeamercolor{title}{fg=white, bg=softorange}
\setbeamercolor{frametitle}{fg=darkgrey}
\setbeamercolor{background canvas}{bg=cream}

% --- PACKAGES ---
\usepackage{amsmath, amsfonts, amssymb}
\usepackage{graphicx}
\usepackage{booktabs}
\usepackage{hyperref}

% --- TITLE PAGE ---
\title[Equity Performance Project]{Equity Performance Project \\ \large Analysis of Investor Personas and Index Tracking}
\author{Juan Pérez Osorio}
\institute{Department of Applied Mathematics and Statistics \\ Stony Brook University}
\date{November 2025}

\begin{document}
	
	% --- TITLE SLIDE ---
	\begin{frame}
		\titlepage
	\end{frame}
	
	% --- MOTIVATION ---
	\begin{frame}{Motivation}
		\begin{itemize}
			\item Modern Portfolio Theory provides a framework for risk–return tradeoffs.
			\item We evaluate a 10-stock equity portfolio under two investor personas.
			\item Risk metrics: Sharpe Ratio, VaR, CVaR, volatility decomposition.
		\end{itemize}
	\end{frame}
	
	% --- PORTFOLIO ---
	\begin{frame}{Portfolio Description}
		\textbf{10-Stock Portfolio:}
		\begin{itemize}
			\item Technology: MSFT, AAPL, IBM
			\item Dividend Aristocrats: XOM, WTM, KO, JNJ
			\item High Dividend Yield: VZ, KHC, USB
		\end{itemize}
	\end{frame}
	
	% --- PERSONAS ---
	\begin{frame}{Investor Personas}
		\textbf{Buy \& Hold (BH):}
		\begin{itemize}
			\item Long-term investor, single wealth trajectory.
			\item P\&L: \( \text{Value}_E - \text{Value}_S \)
		\end{itemize}
		
		\vspace{4mm}
		
		\textbf{Hedge Fund (HF):}
		\begin{itemize}
			\item Daily mark-to-market.
			\item P\&L series: \( \text{Value}_t - \text{Value}_{t-1} \)
		\end{itemize}
	\end{frame}
	
	% --- RISK METRICS ---
	\begin{frame}{Risk and Performance Metrics}
		\begin{itemize}
			\item Sharpe Ratio
			\item Value at Risk (Historical Simulation)
			\item Conditional Value at Risk
			\item Volatility estimation and rolling windows
		\end{itemize}
	\end{frame}
	
	% --- WEIGHTING METHODS ---
	\begin{frame}{Weighting Methodologies}
		\begin{itemize}
			\item Equal Weighting
			\item Market Capitalization Weighting
			\item Logarithmic Market Cap Weighting
		\end{itemize}
	\end{frame}
	
	% --- INDEX TRACKING ---
	\begin{frame}{Index Tracking Optimization}
		\textbf{Objective:} Minimize Tracking Error Variance:
		\[
		\min_w \; \mathbb{E}[(R_T - R_R)^2]
		\]
		\textbf{Constraints:}
		\begin{itemize}
			\item Full investment: \( \sum w_i = 1 \)
			\item Long-only: \( w_i \ge 0 \)
			\item Cardinality: \( \|w\|_0 \le k \)
			\item Buy-in threshold: \( w_i \in \{0\} \cup [l,1] \)
		\end{itemize}
	\end{frame}
	
	% --- RESULTS SUMMARY ---
	\begin{frame}{Summary of Findings}
		\begin{itemize}
			\item Market-cap weighting delivered high returns but extreme concentration.
			\item Equal weighting provided stable diversification but lower risk-adjusted returns.
			\item Log-market-cap weighting offered a balanced alternative.
			\item Index tracking achieved low tracking error under practical constraints.
		\end{itemize}
	\end{frame}
	
	% --- CONCLUSIONS ---
	\begin{frame}{Conclusions}
		\begin{itemize}
			\item Investor objectives strongly influence optimal portfolio design.
			\item Weighting choices drive concentration, turnover, and performance.
			\item Optimization frameworks enable realistic index replication.
		\end{itemize}
	\end{frame}
	
	% --- END ---
	\begin{frame}
		\centering \Large Thank you!
	\end{frame}
	
\end{document}
